\section{Hardware Evaluation}
\label{sec:hweval}

Section~\ref{sec:arch} locked the integer contract, the unified
pipeline, and the session capacity at $N{=}64$. This section evaluates
the resulting IP under a 28-nm CMOS standard-cell flow, establishes the
architectural case for a dedicated scheduling core, and provides a
board-level demonstration. The evaluation proceeds from synthesis
results through efficiency, timeliness, and fidelity analyses to a
structural necessity argument that explains why the same policy cannot
be realized in software or as a pair of independent IPs.

\subsection{Locus Against Datapath KV Hardware}
\label{sec:hw-locus}

Table~\ref{tab:hw-kv} positions the \unison scheduling core against
token-level, device-level, and interconnect-level engines. Every
existing entry operates on a single request and optimizes a per-token
or per-block metric, whereas the scheduler ranks a pool of live
sessions on gap and completion signals that no token-level engine
observes. The two classes of hardware are complementary rather than
competing.

\begin{table*}[t]
\centering
\renewcommand{\arraystretch}{1.12}
\setlength{\tabcolsep}{2pt}
\scriptsize
\begin{threeparttable}
\caption{Comparison of representative KV cache hardware against the
session-level near-memory scheduler.\tnote{a}}
\label{tab:hw-kv}
\begin{tabular*}{\textwidth}{@{\extracolsep{\fill}}lllll@{}}
\toprule
\textbf{Work} & \textbf{Granularity} & \textbf{Locus} &
\textbf{Decision epoch} & \textbf{Mechanism} \\
\midrule
Token-Picker~\cite{park_tokenpicker_2024} & token & attention datapath
  & decode step & lossy probability prune; skip KV transfer \\
UniCAIM~\cite{xu_unicaim_2025} & token & CAM/CIM array
  & attention step & lossy static-dynamic prune; in-place attention \\
KV-MMU~\cite{moradifirouzabadi_kvmmu_2025} & token & accelerator MMU
  & generate step & lossy regularized token replace \\
HiKV~\cite{fang_hikv_2026} & token and element & attention datapath
  & decode step & lossy two-stage importance sort \\
Kelle~\cite{xia_kelle_2025} & token & eDRAM controller
  & decode step & lossy attention-score evict; refresh density \\
CXL-SpecKV~\cite{liu_cxlspeckv_2026} & block & CXL FPGA
  & speculative prefetch & expand pool; predict and move blocks \\
V-Rex~\cite{kim_vrex_2026} & video token & retrieval engine
  & frame / iterative prefill & retrieve a frame subset \\
\rowcolor{gray!10}
\textbf{\textsc{Unison}} & session & near-memory scheduler
  & request and gap events & lossless next-reference rank; evict and tier \\
\bottomrule
\end{tabular*}
\begin{tablenotes}\scriptsize
\item[a] Token-level importance datapaths also include
MATA~\cite{zhu_mata_2026} and VEDA~\cite{wang_veda_2025}.
\end{tablenotes}
\end{threeparttable}
\end{table*}


\subsection{Synthesis Methodology}
\label{sec:hw-method}

The synthesized netlist comprises an APB wrapper, a CSR file, an event
FIFO, and the unified \spear scan plus \tide engine, with the session
table mapped to flip-flops at roughly $N_{\mathrm{sess}}\times 126$
bits plus a four-bit \tide working copy. SRAM compiler macros, the HBM
PHY, and KV payload arrays are outside the cell area, and a prefix
directory together with a multi-instance router are likewise excluded.

Logic synthesis uses Synopsys Design Compiler on a commercial 28-nm
CMOS library at a 150\,MHz clock constraint, an application-side
budget rather than a race against the HBM PHY. Power is a vectorless
estimate at the timing corner. The RTL locks a pipelined scan and the
exact divider of~\eqref{eq:t1}, matching the bit-exact cosimulation of
Section~\ref{sec:micro}. Unless stated otherwise, bit-exact validation
and the FPGA demonstration use the GAIA/Gemma4 Edge-Tight trace, which
is the most hazard-sensitive of the six families and therefore the
hardest case for the fixed-point contract. Timeliness analysis draws on
all three GAIA families because only the GAIA traces carry recorded
tool-gap annotations.

\subsection{Area, Power, and Timing}
\label{sec:hw-ppa}

Table~\ref{tab:dc} reports the 64-session design point at the ASE knee
of Section~\ref{sec:dse}, occupying 0.169\,mm$^{2}$ at 13.6\,mW and
meeting 150\,MHz with 0.01\,ns positive setup slack for an estimated
$F_{\max}$ of 153.4\,MHz.

% Auto-generated by paper/tables/gen_tab_dc.py
\begin{table}[t]
\centering
\renewcommand{\arraystretch}{1.12}
\setlength{\tabcolsep}{3pt}
\scriptsize
\begin{threeparttable}
\caption{Back-end synthesis results of the \unison scheduling core.}
\label{tab:dc}
\begin{tabular*}{\columnwidth}{@{\extracolsep{\fill}}lrrrrr@{}}
\toprule
\textbf{Process} & \textbf{Area}\tnote{a} &
\textbf{Power} & \textbf{WNS} & \textbf{$F_{\max}$} &
\textbf{Cells} \\
 & (mm$^{2}$) & (mW) & (ns) & (MHz) &  \\
\midrule
28-nm CMOS & 0.169 & 13.6 & +0.01 & 153.4 & 158k \\
\bottomrule
\end{tabular*}
\begin{tablenotes}\scriptsize
\item[a] Scope is the
APB wrapper, CSR, event FIFO, and unified core, excluding HBM PHY and KV payload SRAM.
\end{tablenotes}
\end{threeparttable}
\end{table}


Combinational logic accounts for 54\% of cell area and sequential
logic for 46\%. The \tide budget divider is the largest block at 34\%,
followed by the \spear divider at 17\% and the APB wrapper, CSR, and
event FIFO at a combined 6\%. The critical path runs through the first
\tide divider stage from the programmed DMA bandwidth register into the
remainder pipeline.

On the Yosys/Nangate45 flow of Fig.~\ref{fig:dse}, ASE reaches 0.900,
0.952, and 0.985 at 16, 32, and 64 sessions and saturates beyond 64
because trace concurrency never exceeds that count. Mean decision
latency per~\eqref{eq:tau} grows from 0.30\,$\mu$s at 16 sessions to
2.00\,$\mu$s at 64 and 3.95\,$\mu$s at 256. The five structural
variants at $N{=}64$ span 2.6\% in cells, and varying $b_{h}$ or
$b_{s}$ from the locked values shifts area by less than 3\%,
confirming that the precision demanded by ranking is essentially free
relative to the register file.

\subsection{Decision Efficiency and Control-Plane Overhead}
\label{sec:hw-eff}

The scheduler outputs a victim index and DMA descriptors rather than
computing attention, so its relevant efficiency metric is control-plane
cost relative to the data plane it manages.

\paragraph{Area amortization}
The 0.169\,mm$^{2}$ controller manages a two-tier pool of 10\,K to
60\,K tokens. At 128 heads, 128 dimensions, and 16-bit KV per layer, a
60\,K-token HBM bank occupies several hundred megabytes, and the
scheduling core is negligible relative to the HBM PHY and KV payload
arrays that together dominate die area on any contemporary inference
accelerator. Yet this vanishing fraction of total silicon delivers a
34.8\% mean AMAT reduction and a B\'el\'ady ratio of 0.93, a
cost-benefit asymmetry analogous to the branch predictor that occupies
less than 2\% of core area while eliminating the majority of pipeline
stalls.

\paragraph{Energy per decision}
At 13.6\,mW and a 2.00\,$\mu$s mean decision latency, each eviction
or migration decision consumes approximately 27.2\,pJ, which is orders
of magnitude below the per-access energy of the HBM transactions that
the decision orchestrates.

\paragraph{Workload proportionality}
Controller area scales with $N_{\mathrm{sess}}$ through the register
file, not with model dimensionality or sequence length. Increasing model
size from 7\,B to 70\,B parameters multiplies the KV store and
attention datapath but leaves the scheduling core unchanged, while
doubling session capacity beyond 64 roughly doubles the register file
without improving ASE, which is already saturated at 0.985. A
general-purpose microcontroller could in principle execute the same
algorithm, but its instruction-fetch and interrupt overhead would push
decision latency into the tens of microseconds, and the area of even a
minimal core with tightly coupled memory exceeds the 0.169\,mm$^{2}$
of the dedicated datapath while delivering lower determinism.

\subsection{Decision Timeliness}
\label{sec:hw-time}

At $N{=}64$ and 150\,MHz the mean scan latency is 2.00\,$\mu$s and
the worst case is 3.14\,$\mu$s. The GAIA traces exhibit median tool
gaps of 3.7\,s for Qwen3, 4.4\,s for Devstral, and 7.2\,s for
Gemma4, with the minimum observed gap across all three families at
1\,ms. The worst-case decision latency is therefore 0.31\% of the
minimum gap and below 0.0001\% of the median, leaving four to six
orders of magnitude of timing headroom.

This margin implies that even a ten-fold frequency reduction to
15\,MHz would keep the decision latency at 31\,$\mu$s, well within the
millisecond-scale minimum gap, so the controller can be placed in a
slow clock domain or power-gated between events. By contrast, a
software scheduler polling block-cache timestamps at millisecond
granularity occupies the same order of magnitude as the minimum gap
itself, which is why the software approximation inverts the gap signal
on short waits as documented in Section~\ref{sec:sw-vllm}.

\subsection{Implementation Fidelity}
\label{sec:hw-fid}

Cosimulation of the locked RTL format against the floating-point
reference is bit-exact on all 1\,937 events and 845 decisions of the
validation trace. Across 7\,265 real eviction snapshots from six traces and
three capacity envelopes, the aggregate top-1 victim agreement is
0.9945, the Kendall $\tau$ exceeds 0.998, and per-dataset agreement
ranges from 0.991 on SWE/Devstral at 2\,400 snapshots to 1.000 on
GAIA/Devstral at 16 snapshots, with no dataset falling below the 0.99
hard gate. The worst-cell end-to-end hit-rate drift is 0.192\,pp
within the 0.20\,pp hard gate, and the mean drift across all 18
dataset-by-capacity cells is 0.019\,pp while the AMAT drift is
$-0.018$\%.

Because each eviction decision is an independent $\arg\max$ over
instantaneous scores, the 0.55\% top-1 disagreement is a per-snapshot
coin flip that does not accumulate into trajectory-level drift, as
confirmed by the 0.019\,pp mean hit-rate deviation over full traces.
The fixed-point pipeline is therefore functionally transparent with
respect to the policy it implements, reproducing the ranking within a
residual that is smaller than the inter-policy gap between \unison and
CacheScout on every dataset.

\subsection{Structural Necessity of the Unified Near-Memory Design}
\label{sec:hw-nec}

Three lines of evidence converge to show that the same policy cannot be
realized more cheaply in software or as a pair of independent IPs.

\paragraph{Software cannot reconstruct the event interface}
Section~\ref{sec:sw-vllm} demonstrates that a software scheduler
operating on block-cache timestamps inverts the gap signal when the
tool wait approaches the decode window, because the timestamp
granularity at millisecond resolution is coarser than the gap itself,
whereas the hardware event FIFO resolves the same transition at cycle
granularity. More broadly, any software scheduler sharing the Python
runtime is subject to GIL serialization and OS-level jitter that
inflate scheduling latency by orders of magnitude beyond the
microsecond decisions the hardware achieves.

\paragraph{A dual-IP design diverges under agent interleaving}
Section~\ref{sec:org} shows that splitting eviction and migration into
two engines with private register files causes hit-rate drops of up to
5.4\,pp even under every-event synchronization, because agent workloads
interleave request and gap events at fine granularity rather than in
long homogeneous phases.

\paragraph{Convergence of the three constraints}
Cycle-level event observation requires near-memory placement, consistent
state requires a unified register file, and deterministic gap-window
completion requires a pipelined datapath. Relaxing any one constraint
re-introduces a documented failure mode, so the architectural
contribution is this constraint intersection realized as a
sub-0.2\,mm$^{2}$ control-plane IP.

\subsection{Board-Level Demonstration}
\label{sec:hw-fpga}

Fig.~\ref{fig:fpga} shows a Xilinx Zynq-7020 prototype that replays
the validation trace over 1\,937 events and 845 decisions, stamping
decision latency at the 150\,MHz Design Compiler clock. The
post-implementation netlist occupies 43\,415 LUTs at 81.6\% slice
utilization, 20\,938 flip-flops at 19.7\%, and 20 DSP48E1 blocks at
9.1\%, with no block RAM consumed.

\begin{figure}[t]
\centering
\includegraphics[width=\columnwidth]{fpga_demo_annotated.jpg}
\caption{FPGA prototype implementation of the \unison scheduling core
on Xilinx Zynq-7020.}
\label{fig:fpga}
\end{figure}

\subsection{Physical Placement Considerations}
\label{sec:hw-place}

The timeliness margin of Section~\ref{sec:hw-time} implies that the
controller does not require co-location with the HBM PHY clock domain.
The essential requirement is access to the event stream that software
timestamps cannot reconstruct, which a sideband event bus from the
memory controller provides without a shared clock or a high-bandwidth
data interface.
